\documentclass[10pt,a4paper]{amsart}
\usepackage{amsmath,amssymb,amsthm,graphicx}
% For commutative diagrams we use Xy-pic
\usepackage[all]{xy}

% We use 2cell for 2-commutative diagrams.
\xyoption{2cell}
% Theorem environments.
%
\theoremstyle{plain}
\newtheorem{theorem}[subsection]{Theorem}
\newtheorem{proposition}[subsection]{Proposition}
\newtheorem{lemma}[subsection]{Lemma}

\theoremstyle{definition}
\newtheorem{definition}[subsection]{Definition}
\newtheorem{example}[subsection]{Example}
\newtheorem{exercise}[subsection]{Exercise}
\newtheorem{situation}[subsection]{Situation}

\theoremstyle{remark}
\newtheorem{remark}[subsection]{Remark}
\newtheorem{remarks}[subsection]{Remarks}

\numberwithin{equation}{subsection}

% Macros
%
\def\lim{\mathop{\mathrm{lim}}\nolimits}
\def\colim{\mathop{\mathrm{colim}}\nolimits}
\def\Spec{\mathop{\mathrm{Spec}}}
\def\Hom{\mathop{\mathrm{Hom}}\nolimits}
\def\Ext{\mathop{\mathrm{Ext}}\nolimits}
\def\SheafHom{\mathop{\mathcal{H}\!\mathit{om}}\nolimits}
\def\SheafExt{\mathop{\mathcal{E}\!\mathit{xt}}\nolimits}
\def\Sch{\mathit{Sch}}
\def\Mor{\mathop{\mathrm{Mor}}\nolimits}
\def\Ob{\mathop{\mathrm{Ob}}\nolimits}
\def\Sh{\mathop{\mathit{Sh}}\nolimits}
\def\NL{\mathop{N\!L}\nolimits}
\def\CH{\mathop{\mathrm{CH}}\nolimits}
\def\proetale{{pro\text{-}\acute{e}tale}}
\def\etale{{\acute{e}tale}}
\def\QCoh{\mathit{QCoh}}
\def\Ker{\mathop{\mathrm{Ker}}}
\def\Im{\mathop{\mathrm{Im}}}
\def\Coker{\mathop{\mathrm{Coker}}}
\def\Coim{\mathop{\mathrm{Coim}}}

% Boxtimes
%
\DeclareMathSymbol{\boxtimes}{\mathbin}{AMSa}{"02}

%
% Macros for moduli stacks/spaces
%
\def\QCohstack{\mathcal{QC}\!\mathit{oh}}
\def\Cohstack{\mathcal{C}\!\mathit{oh}}
\def\Spacesstack{\mathcal{S}\!\mathit{paces}}
\def\Quotfunctor{\mathrm{Quot}}
\def\Hilbfunctor{\mathrm{Hilb}}
\def\Curvesstack{\mathcal{C}\!\mathit{urves}}
\def\Polarizedstack{\mathcal{P}\!\mathit{olarized}}
\def\Complexesstack{\mathcal{C}\!\mathit{omplexes}}
% \Pic is the operator that assigns to X its picard group, usage \Pic(X)
% \Picardstack_{X/B} denotes the Picard stack of X over B
% \Picardfunctor_{X/B} denotes the Picard functor of X over B
\def\Pic{\mathop{\mathrm{Pic}}\nolimits}
\def\Picardstack{\mathcal{P}\!\mathit{ic}}
\def\Picardfunctor{\mathrm{Pic}}
\def\Deformationcategory{\mathcal{D}\!\mathit{ef}}
\newenvironment{reference}{}{}
\newenvironment{slogan}{}{}
\newenvironment{history}{}{}
\title{The recurring theorem on determinants: diagonal dominance with arbitrary positive weights forces a nonzero determinant}
\author{The Stacks Project Authors}
\date{stacks-project:models.tex:local-snapshot-20261012; GNU FDL 1.2 or later}
\begin{document}
\maketitle
\noindent\textit{Editorial note.} The counted unit is the Stacks Project Lemma lemma-recurring at source lines 303-316 together with its proof at source lines 318-328, a straight adjacent pairing in the chapter file. The chapter is the upstream file models.tex, which is the file name recorded in the source record; the pool directory carries the longer name stacks-models and is not used as the chapter name anywhere in this card.\par\medskip
\noindent\textit{Editorial note.} This chapter's candidate index records the statement body and the proof body without their closing markers, so the retained spans used here are the full environments: the statement span starts at the opening marker of the lemma environment and ends at its closing marker, and the proof span likewise runs from the opening to the closing marker of the proof environment. The arithmetic-progression candidate index of this chapter was built with the same convention as the earlier cards of this chapter.\par\medskip
\noindent\textit{Editorial note.} The statement begins with the author's own reference environment, which is a comment environment in this chapter and contains the citation of the quoted theorem. The retained preamble keeps the wrapper that makes the reference environment visible for this lane, so the citation is inside visible material rather than discarded, and one bibliography entry is supplied as a task-directory fragment holding it exactly as a bibtex run on the host had formatted it out of the Stacks Project's shared bibliography. The bibliography backend is therefore none and no bibtex pass is performed.\par\medskip
\noindent\textit{Editorial note.} The quoted citation carries an optional argument naming a numbered theorem of the cited paper, and that optional argument is retained verbatim in the span.\par\medskip
\noindent\textit{Editorial note.} The retained preamble is the chapter preamble's Xy-pic setup at lines 14-18, the theorem declarations and mathematical macros at lines 154-217, and the three no-op wrapper environments.\par\medskip
\noindent\textit{Editorial note.} Neither the retained statement nor the retained proof refers to another labelled result of this or of any other chapter, so this card carries no navigation pointer.\par\medskip
\noindent\textit{Allowed standard prerequisites (editorial).}\begin{itemize}
\item The quoted determinant theorem and the notion of a recurring matrix. The reference environment at the top of the counted statement quotes the classical theorem on determinants of a matrix whose diagonal entries dominate their rows, in the two forms that the counted statement then re-proves for complex matrices and for the analogue obtained by rescaling the columns with a positive real vector. The chapter's own name for the second form is that of a recurring matrix.. The counted lemma is the chapter's own write-up of that quoted theorem, and the surrounding section uses it to show that the intersection matrix of a numerical type has rank one less than its size and that its kernel is spanned by the weight vector.
\end{itemize}\medskip
\begin{lemma}
\label{lemma-recurring}
\begin{reference}
\cite[Theorem I]{Taussky}
\end{reference}
Let $A = (a_{ij})$ be a complex $n \times n$ matrix.
\begin{enumerate}
\item If $|a_{ii}| > \sum_{j \not = i} |a_{ij}|$ for each $i$, then
$\det(A)$ is nonzero.
\item If there exists a real vector $m = (m_1, \ldots, m_n)$
with $m_i > 0$ such that $|a_{ii} m_i| > \sum_{j \not = i} |a_{ij}m_j|$
for each $i$, then $\det(A)$ is nonzero.
\end{enumerate}
\end{lemma}

\begin{proof}
If $A$ is as in (1) and $\det(A) = 0$, then there is a nonzero vector
$z$ with $Az = 0$. Choose $r$ with $|z_r|$ maximal. Then
$$
|a_{rr} z_r| = |\sum\nolimits_{k \not = r} a_{rk}z_k| \leq
\sum\nolimits_{k \not = r} |a_{rk}||z_k| \leq
|z_r| \sum\nolimits_{k \not = r} |a_{rk}| < |a_{rr}||z_r|
$$
which is a contradiction. To prove (2) apply (1) to the matrix
$(a_{ij}m_j)$ whose determinant is $m_1 \ldots m_n \det(A)$.
\end{proof}

\begin{thebibliography}{10}

\bibitem{Taussky}
O.~Taussky.
\newblock A recurring theorem on determinants.
\newblock {\em Amer. Math. Monthly}, 56:672--676, 1949.

\end{thebibliography}
\end{document}
